\scanpage{13}
$M(S)$ définissent des courbes elliptiques isomorphes (\unsure{non|bien} structurées) il f et s que \deleted{\unsure{elles}} ils soient congrus par une opération de $G(S)$, où \inserted{$G$ est un groupe simple sur $S$.} Enfin, pour tte courbe elliptique $C$ sur $S$, il existe une extension \deleted{fidt plate} \inserted{simple surj. de type fini} $S'$ de $S$ [savoir le fibré principal associé à $\mathcal E_3$ pour le clivage qu'on a mis sur $\mathcal E_3$] tel que \deleted{\unsure{sur}} $C\times_S S'$ ``provienne'' d'un élément de $M(S')$.

Ainsi, avec une terminologie qui court, le foncteur $F$ : ``classes à isom près de courbes elliptiques sur $S$'' est \deleted{\unsure{fppf}}-$\mathcal D$-représentable, où $D$ désigne la catégorie des morphismes simples surjectifs. Si ce n'était que pour obtenir celà d'ailleurs on n'aurait pas à se donner tant de mal, \uline{car} \unsure{là|de} on a un raisonnement général qui le donne pour les courbes de genre $g$ \unsure{quelc} (le groupe triangulaire remplacé par un groupe projectif ---).
% archive sheet number: 12
